\section{Identification: what each plot needs from the data}
\label{sec:identification}

\emph{Source: handout \S1.1. This is the intern-term paper's (``Paper A'')
first result. Status: skeleton from a first read.}

\subsection{The reduced form}

Let $R$ be the features other than $A$ and its parents $W$. Let
$Q(dr \mid a, w)$ be the conditional law of $R$ given $(A,W)$ --- the
\textbf{reduced form}. Crucially, this is \emph{learnable straight from
data} --- no ECM structural equations needed.

\subsection{Per-plot identification table}

\begin{table}[h]
\centering
\begin{tabular}{@{}p{2.1cm}p{4.4cm}p{3.6cm}p{3.4cm}@{}}
\toprule
\textbf{Plot} & \textbf{Population target} & \textbf{What it consumes} &
\textbf{Statistical character} \\
\midrule
TDP & $\int P(dw) \int Q(dr\mid a,w)\, f(a,w,r)$ & parents $+$ reduced form &
  AIPW; $\sqrt{n}$ at discrete levels \\
NDDP $=$ PDP & $E[f(a, X_{-A})]$ & \textbf{no causal nuisance at all} ---
  only $f$ and the reference distribution of $X_{-A}$ & sample mean of a
  known function; $\sqrt{n}$ at each $a$ \\
NIDP & $\int P(da', dw) \int Q(dr\mid a,w)\, f(a', w, r)$ & parents $+$
  reduced form $+$ law of $A$ given $W$ & its own influence function;
  $\sqrt{n}$ at discrete levels \\
PCDP & needs the graph \emph{inside} $R$ & more than the reduced form (two
  graphs with the same data give different answers) & g-formula \\
\bottomrule
\end{tabular}
\caption{Identification of each named plot (handout \S1.1).}
\label{tab:identification}
\end{table}

\emph{Reading the three dense formulas in words:} TDP's
$\int P(dw) \int Q(dr\mid a,w)\, f(a,w,r)$ says --- go age group by age
group ($\int P(dw)$, weighting by how common each age group actually is);
within each age group, average over what weight-and-blood-pressure values
typically go with exercise level $a$ there ($\int Q(dr\mid a,w)$, the
reduced form); then plug each $(a,w,r)$ combination into the model $f$ and
average. NIDP's formula is almost identical but starts by also averaging
over exercise levels $a'$ drawn the way they actually occur for each age
group ($\int P(da',dw)$) before asking what downstream values go with the
\emph{target} level $a$ --- that's the ``your own exercise, but their
downstream values'' mixing. NDDP's $E[f(a,X_{-A})]$ is the simplest: just
plug the fixed value $a$ in for exercise, leave every other feature exactly
as recorded in the data, run it through $f$, and average --- no age
groups, no reduced form, nothing learned about causal structure at all.

\textbf{Takeaway to hold onto:} three of the four named plots (TDP, NDDP,
NIDP) need \emph{only} the parents of $A$ plus one reduced form --- fitting
every equation of the ECM is never required. Only PCDP needs the internal
graph structure of $R$. This is what makes ``adjustment''
(\S\ref{sec:estimation}) a real alternative to ``propagation''
(\S\ref{sec:propagation}) for three of the four plots.

\subsection{NIDP \texorpdfstring{$=$}{=} a known bridge mean (reuse, not a new discovery)}

\reusetag{Fulcher, Shpitser, Marealle \& Tchetgen Tchetgen 2020, JRSSB}
``Your own $A$, but the downstream values you'd have had under $A=a$'' is
exactly their \textbf{bridge mean} $\Psi(a) = E[Y(A, Z(a))]$ from the
population-intervention indirect-effect literature ($Z=$ downstream
features, $Y=f$) --- in words, plug each person's \emph{own actual}
exercise level ($A$, not $a$) into the model, but pair it with the
downstream feature values ($Z(a)$) they \emph{would have had} under the
hypothetical exercise level $a$, then average over everyone. It mixes one
real, observed input with one counterfactual input for the same person.
Their Lemma~1 gives the formula, their Theorem~1 gives its
influence function (one term drops out here because $Y^*=f(X)$ is
\emph{deterministic}, not noisy), and their doubly-robust estimator applies
directly.

What's \newtag~here specifically: recognizing the published NIDP \emph{is}
this object, that the covariate set reduces to just the parents of $A$, and
the simplifications that follow from $f$ being a known function rather than
an unknown outcome-generating process. Caveat: handled for
\textbf{discrete $A$ only} this term --- continuous-$A$ NIDP needs its own
theory.

\subsection{Direct \texorpdfstring{$+$}{+} indirect \texorpdfstring{$\neq$}{!=} total}

\newtag, modest. With $B$ the average prediction:
\[
  \theta(a) - B \;=\; [\mathrm{NDDP}(a) - B] + [\mathrm{NIDP}(a) - B] + J(a),
\]
(in words: the total curve's distance from the average, on the left, isn't
simply the direct effect's distance plus the indirect effect's distance ---
there's a leftover term $J(a)$ needed to make the two sides balance; see
\S\ref{sec:bigpicture} for the full walkthrough). $J$ is an interaction remainder --- it vanishes only when $f$ is additive in
$a$ and the downstream features. Toy counterexample: $A \sim \mathrm{Bern}(p)$,
$R(a)=a$, $f = a\cdot r$: the total contrast is $1$, but the direct and
indirect contrasts are each only $p$. \textbf{Don't assume direct $+$
indirect sums to total} --- $J$ should be plotted alongside them, not
ignored.

\begin{openquestions}
  \item Why does PCDP specifically need the graph \emph{inside} $R$, when
    TDP/NDDP/NIDP don't need any equations at all? (Intuition: fixing an
    intermediate feature ``cuts'' a path in a way sensitive to which
    downstream variable causes which --- worth confirming with a toy SCM.)
  \item Verify \S1.1's table by exact enumeration in toy SCMs --- handout
    \S1.9 marks this as intern-feasible / light supervision (\feasyes).
    Good first task.
\end{openquestions}
